% =====================================================================
% Reviewer 1, comment 7 -- rewritten Section 3.1.3 (Results, quantitative).
% Numbers come from validation/r1c7_ablation.py (17 sites, 1991-2020).
% Renamed to signal a variable-importance result, not a qualitative remark.
% =====================================================================
\subsubsection{Variable importance: solar radiation and climatological priors}
\label{sec:ablation}

To quantify how much each pipeline component contributes to the final skill, we
performed a component-isolation (ablation) experiment on the 17~MATOPIBA and
Piracicaba sites (1991--2020; \Nobs~daily observations). Starting from the full
EVAonline configuration, we disabled \emph{one} component at a time and measured
the change in KGE and percent bias (PBIAS) against BR-DWGD
(Table~\ref{tab:ablation}). Because only one element changes per row, each
$\Delta$KGE is directly attributable to that element.

\begin{table}[!htbp]
\centering
\caption{Ablation experiment: change in performance when a single EVAonline
component is removed (17 sites, 1991--2020). Negative $\Delta$KGE indicates loss
of skill; PBIAS is the site-wise mean.}
\label{tab:ablation}
\small
\begin{tabular}{lcccc}
\toprule
Configuration & KGE & PBIAS (\%) & $\Delta$KGE & $\Delta$PBIAS \\
\midrule
Baseline (full EVAonline) & 0.812 & +0.36 & -- & -- \\
A1: $R_s$ from Open-Meteo only (no CERES weight) & 0.756 & +0.51 & $-0.056$ & $+0.15$ \\
A2: $R_s$ replaced by climatology & 0.637 & $-0.41$ & $-0.175$ & $-0.77$ \\
A3: no climatological bias correction & 0.722 & $+10.87$ & $-0.091$ & $+10.50$ \\
A4: no adaptive outlier rejection & 0.809 & $-0.00$ & $-0.003$ & $-0.36$ \\
A5: no dynamic $Q$ adaptation & 0.706 & $+0.22$ & $-0.106$ & $-0.15$ \\
\bottomrule
\end{tabular}
\end{table}

The experiment confirms that \textbf{solar radiation is the single most
influential input}: replacing the daily $R_s$ with its monthly climatology (A2)
removes most of the day-to-day signal and is the largest individual degradation
($\Delta$KGE${}=-0.175$), while down-weighting the CERES-derived radiation in
Stage~1 (A1) costs a further $\Delta$KGE${}=-0.056$. This agrees with the
tropical sensitivity analyses of \citet{Irizarry2023} and \citet{Paca2022} and
justifies prioritising CERES radiation in the fusion weights (Table~6).

The climatological priors contribute through two distinct mechanisms. The
monthly bias correction (A3) is what keeps the estimates unbiased: disabling it
leaves KGE comparatively high (0.722) but inflates PBIAS from $+0.36\%$ to
$+10.9\%$, i.e.\ the priors are responsible for the near-zero bias that
distinguishes EVAonline from the single-source baselines
(Table~\ref{tab:validation-results}). The dynamic process-noise adaptation (A5,
$\Delta$KGE${}=-0.106$) lets the filter follow genuine sharp transitions, whereas
the asymmetric p01/p99 outlier rejection (A4) has a negligible \emph{aggregate}
effect ($\Delta$KGE${}=-0.003$) because it acts only on a small number of
extreme days; its benefit is therefore concentrated in the tails rather than in
the 30-year aggregate statistics. Together, these results show that EVAonline's
skill arises from a physically meaningful division of labour: radiation fidelity
sets the daily dynamics, while the climatological priors remove bias and
stabilise the estimate against anomalous inputs.
